% !TEX program = pdflatex
\documentclass[10pt]{article}

\usepackage[utf8]{inputenc}
\usepackage[T1]{fontenc}
\usepackage[english]{babel}
\usepackage{lmodern}
\usepackage{microtype}
\usepackage[margin=0.72in]{geometry}
\usepackage{amsmath,amssymb,amsthm,mathtools}
\usepackage{booktabs}
\usepackage{xcolor}
\usepackage{enumitem}
\usepackage[most]{tcolorbox}
\usepackage{fancyhdr}
\usepackage[hidelinks]{hyperref}

\definecolor{CourseBlue}{HTML}{1A3C68}
\definecolor{CourseGreen}{HTML}{216E39}
\definecolor{CourseGold}{HTML}{B8860B}
\definecolor{CourseRed}{HTML}{9A2C2C}
\definecolor{SoftBlue}{HTML}{F2F6FA}
\definecolor{SoftGreen}{HTML}{F1F8F3}
\definecolor{SoftGold}{HTML}{FFF8E8}
\definecolor{SoftRed}{HTML}{FFF3F3}

\setlength{\parindent}{0pt}
\setlength{\parskip}{0.32em}
\setlist[itemize]{leftmargin=1.5em,itemsep=0.12em,topsep=0.2em}
\setlist[enumerate]{leftmargin=1.7em,itemsep=0.12em,topsep=0.2em}
\allowdisplaybreaks

\pagestyle{fancy}
\fancyhf{}
\lhead{Stochastic Calculus and Option Pricing}
\rhead{Lecture 3 Notes}
\cfoot{\thepage}
\renewcommand{\headrulewidth}{0.3pt}

\newcommand{\E}{\mathbb E}
\newcommand{\Pbb}{\mathbb P}
\newcommand{\R}{\mathbb R}
\newcommand{\1}{\mathbf 1}
\newcommand{\cF}{\mathcal F}
\newcommand{\dd}{\,\mathrm d}
\newcommand{\Ito}{\mathrm{It\hat{o}}}
\newcommand{\norm}[1]{\left\lVert #1\right\rVert}
\DeclareMathOperator{\Var}{Var}
\DeclareMathOperator{\Cov}{Cov}

\theoremstyle{definition}
\newtheorem{definition}{Definition}[section]
\newtheorem{proposition}[definition]{Proposition}
\newtheorem{theorem}[definition]{Theorem}
\newtheorem{corollary}[definition]{Corollary}
\newtheorem{remark}[definition]{Remark}

\newtcolorbox{boardexample}[1]{
  enhanced,breakable,colback=SoftBlue,colframe=CourseBlue,
  boxrule=0.55pt,arc=1mm,left=1.5mm,right=1.5mm,top=1mm,bottom=1mm,
  before skip=0.55em,after skip=0.55em,
  title={Example --- #1},fonttitle=\bfseries}
\newtcolorbox{studentexercise}[1]{
  enhanced,breakable,colback=SoftGold,colframe=CourseGold,
  boxrule=0.55pt,arc=1mm,left=1.5mm,right=1.5mm,top=1mm,bottom=1mm,
  before skip=0.55em,after skip=0.35em,
  title={Exercise --- #1},fonttitle=\bfseries}
\newtcolorbox{shortsolution}[1][]{
  enhanced,breakable,colback=SoftGreen,colframe=CourseGreen,
  boxrule=0.55pt,arc=1mm,left=1.5mm,right=1.5mm,top=1mm,bottom=1mm,
  before skip=0.25em,after skip=0.6em,
  title={Short solution#1},fonttitle=\bfseries}
\newtcolorbox{keyidea}[1]{
  enhanced,breakable,colback=white,colframe=CourseGreen,
  boxrule=0.7pt,arc=1mm,left=1.5mm,right=1.5mm,top=1mm,bottom=1mm,
  before skip=0.5em,after skip=0.5em,
  title={#1},fonttitle=\bfseries}
\newtcolorbox{warningbox}[1]{
  enhanced,breakable,colback=SoftRed,colframe=CourseRed,
  boxrule=0.65pt,arc=1mm,left=1.5mm,right=1.5mm,top=1mm,bottom=1mm,
  before skip=0.5em,after skip=0.5em,
  title={#1},fonttitle=\bfseries}

\title{\textbf{Stochastic Integration: Intuition, Construction, and Applications}\\
\large Lecture 3 Notes --- Class Support and Student Summary}
\author{}
\date{}

\begin{document}
\maketitle
\vspace{-1.4em}

\begin{keyidea}{Goals and logical structure}
The stochastic integral is built from left-point trading gains. The construction
has three layers:
\begin{enumerate}
  \item define the integral for simple predictable processes;
  \item prove the \(\Ito\) isometry;
  \item extend by completion in \(L^2(\Omega\times[0,T])\).
\end{enumerate}
Under the expected finite-energy condition, the resulting gain process is a
square-integrable martingale. Deterministic integrands then give explicit
Gaussian applications, while the example \(\int W\dd W\) reveals the role of
Brownian quadratic variation. The final section translates the construction
into reliable simulation rules.
\end{keyidea}

% ============================================================
\section{Why Brownian Motion Needs a New Integral}
% ============================================================

\subsection{From discrete trading gains to a continuous-time target}

Fix a grid
\[
0=t_0<t_1<\cdots<t_n=T.
\]
If \(\phi_{t_i}\) is the number of shares selected at time \(t_i\) and held
over \((t_i,t_{i+1}]\), the self-financing gain is
\[
G_T^\pi=\sum_{i=0}^{n-1}\phi_{t_i}
\bigl(S_{t_{i+1}}-S_{t_i}\bigr).
\]
The continuous-time notation
\[
\int_0^T\phi_t\dd S_t
\]
should preserve this economic meaning: position chosen before a price move,
multiplied by that fresh move, and then summed through time.

For a differentiable path \(S\), ordinary calculus would give
\[
\int_0^T\phi_t\dd S_t=\int_0^T\phi_tS'_t\dd t.
\]
More generally, Riemann--Stieltjes integration works well when the integrator
has finite variation. Brownian paths have neither property: almost surely they
are nowhere differentiable and have infinite total variation on every
nontrivial interval.

The scaling already explains the difficulty. Over an interval of length
\(\Delta t\),
\[
\Delta W\sim\mathcal N(0,\Delta t),
\qquad |\Delta W|\text{ is typically of order }\sqrt{\Delta t}.
\]
On a uniform grid with \(n\) intervals, the sum of absolute increments is
therefore heuristically of order
\[
n\sqrt{T/n}=\sqrt{nT}\longrightarrow\infty.
\]

\subsection{Quadratic variation is the surviving scale}

For a partition \(\pi\), define
\[
Q_T^\pi(W)=\sum_i\bigl(W_{t_{i+1}}-W_{t_i}\bigr)^2.
\]

\begin{theorem}[Quadratic variation on a uniform grid]
If \(t_i=iT/n\), then
\[
Q_T^{(n)}(W)\xrightarrow[n\to\infty]{L^2}T.
\]
In particular, the convergence holds in probability.
\end{theorem}

\begin{proof}
Write \(\Delta W_i=W_{t_{i+1}}-W_{t_i}\) and \(\Delta t=T/n\). The
increments are independent and \(\Delta W_i\sim\mathcal N(0,\Delta t)\).
Thus
\[
\E[Q_T^{(n)}]=\sum_i\E[(\Delta W_i)^2]=n\Delta t=T.
\]
For a centered Gaussian variable with variance \(\sigma^2\),
\(\Var(Z^2)=2\sigma^4\). Hence
\[
\Var(Q_T^{(n)})
=\sum_i\Var((\Delta W_i)^2)
=2n(\Delta t)^2
=\frac{2T^2}{n}\longrightarrow0.
\]
Therefore \(\E[(Q_T^{(n)}-T)^2]\to0\).
\end{proof}

The key point is the change of scale:
\[
\sum_i|\Delta W_i|\longrightarrow\infty,
\qquad
\sum_i(\Delta W_i)^2\longrightarrow T.
\]
The second-order term does not disappear under refinement. It will create the
correction terms characteristic of stochastic calculus.

\subsection{The sampling convention matters}

Take the random integrand \(\phi_t=W_t\) and compare the left- and right-point
sums
\[
L^\pi=\sum_iW_{t_i}\Delta W_i,
\qquad
R^\pi=\sum_iW_{t_{i+1}}\Delta W_i.
\]
Since \(W_{t_{i+1}}=W_{t_i}+\Delta W_i\),
\[
R^\pi-L^\pi=\sum_i(\Delta W_i)^2\longrightarrow T.
\]
Thus the two natural-looking conventions differ by an order-one quantity.
There is no convention-free object called \(\int W\dd W\). The \(\Ito\)
integral chooses left endpoints.

\subsection{Filtration and predictability}

Work on a filtered probability space
\((\Omega,\cF,(\cF_t)_{t\geq0},\Pbb)\) supporting Brownian motion \(W\).
At time \(t_i\), the trader knows \(\cF_{t_i}\). A legal position for the
next interval must be \(\cF_{t_i}\)-measurable.

\begin{definition}[Discrete predictability]
The coefficient used on \((t_i,t_{i+1}]\) is predictable when it is determined
by the information available at \(t_i\). Equivalently, it may depend on the
past and present, but not on \(\Delta W_i=W_{t_{i+1}}-W_{t_i}\).
\end{definition}

This is why left endpoints are not a cosmetic convention. They encode the
information constraint of non-anticipative trading.

\begin{studentexercise}{which rules are admissible?}
On the grid \(0,1,2\), decide whether each rule is predictable:
\begin{enumerate}
  \item \(\phi_0=1\), \(\phi_1=1+W_1^2\);
  \item \(\phi_0=\operatorname{sign}(W_1)\), \(\phi_1=0\);
  \item \(\phi_0=0\), \(\phi_1=\1_{\{W_1>0\}}\);
  \item \(\phi_i=\operatorname{sign}(\Delta W_i)\).
\end{enumerate}
\end{studentexercise}

\begin{shortsolution}
Rules 1 and 3 are admissible. In rule 1, \(1+W_1^2\) is known at time 1;
in rule 3, \(\{W_1>0\}\in\cF_1\). Rule 2 uses \(W_1\) at time 0, and rule
4 uses the future increment itself, so both anticipate.
\end{shortsolution}

If \(\phi_{t_i}\in\cF_{t_i}\) and the products are integrable, then
\[
\E[\phi_{t_i}\Delta W_i\mid\cF_{t_i}]
=\phi_{t_i}\E[\Delta W_i\mid\cF_{t_i}]=0.
\]
Consequently, \(\E[G_T^\pi]=0\). A predictable strategy can reshape risk,
but cannot manufacture an average drift from Brownian increments.

% ============================================================
\section{Construction from Simple Predictable Processes}
% ============================================================

\subsection{The integrability class}

The intended integrands satisfy two conditions:
\[
\phi\text{ is predictable},
\qquad
\E\int_0^T\phi_t^2\dd t<\infty.
\]
The second assumption is called the \emph{expected finite-energy condition}.
It means precisely that
\[
\phi\in L^2(\Omega\times[0,T],\Pbb\otimes\dd t).
\]
The square is natural because a Brownian increment over length \(\Delta t\)
has variance \(\Delta t\): the contribution of a position \(\phi\) to gain
variance is therefore proportional to \(\phi^2\Delta t\).

\begin{warningbox}{Integrability is part of the definition, not an afterthought}
Predictability controls information; square integrability controls size. Both
are needed for the \(L^2\) construction. Later, the same assumption will imply
that the gain at every time belongs to \(L^2\subset L^1\), which is essential
before one may call the gain process a martingale.
\end{warningbox}

\subsection{Definition for simple processes}

\begin{definition}[Simple predictable process]
Fix \(0=t_0<t_1<\cdots<t_n=T\). A simple predictable process has the form
\[
\phi_t=\sum_{i=0}^{n-1}\xi_i\1_{(t_i,t_{i+1}]}(t),
\qquad
\xi_i\in\cF_{t_i},
\]
with \(\E[\xi_i^2]<\infty\). The coefficient \(\xi_i\) is selected at time
\(t_i\) and held over the next interval.
\end{definition}

\begin{definition}[\(\Ito\) integral of a simple process]
For such a process, define
\[
\boxed{
\int_0^T\phi_t\dd W_t
:=\sum_{i=0}^{n-1}\xi_i
\bigl(W_{t_{i+1}}-W_{t_i}\bigr).}
\]
\end{definition}

The definition is literally the discrete gain formula. It immediately gives
the sanity checks
\[
\int_0^Tc\dd W_t=cW_T,
\qquad
\int_0^T\1_{(a,b]}(t)\dd W_t=W_b-W_a,
\]
and time additivity
\[
\int_s^t\phi_r\dd W_r
=\int_s^u\phi_r\dd W_r+\int_u^t\phi_r\dd W_r,
\qquad s\leq u\leq t.
\]
The map is linear in the integrand. Processes that agree
\((\Pbb\otimes\dd t)\)-almost everywhere have the same integral almost surely.

\begin{boardexample}{a deterministic step strategy}
Let \(T=3\) and
\[
\phi_t=2\1_{(0,1]}(t)-\1_{(1,3]}(t).
\]
Then
\[
I=\int_0^3\phi_t\dd W_t
=2W_1-(W_3-W_1).
\]
The two increments are independent, so
\[
I\sim\mathcal N(0,4\cdot1+1\cdot2)=\mathcal N(0,6).
\]
Each interval contributes
\((\text{position})^2\times(\text{interval length})\) to the variance.
\end{boardexample}

\begin{studentexercise}{switch after observing \(W_1\)}
Let \(A=\{W_1>0\}\) and
\[
\phi_t=\1_A\1_{(1,2]}(t).
\]
Explain why \(\phi\) is predictable, write its integral without integral
notation, compute its first two moments, and determine whether its law is
Gaussian.
\end{studentexercise}

\begin{shortsolution}
Since \(A\in\cF_1\), the position is known before \((1,2]\), and
\[
I=\int_0^2\phi_t\dd W_t=\1_A(W_2-W_1).
\]
The event \(A\) is independent of the fresh increment
\(W_2-W_1\sim\mathcal N(0,1)\). Hence
\[
\E[I]=0,
\qquad
\E[I^2]=\Pbb(A)\E[(W_2-W_1)^2]=\frac12.
\]
However, \(\Pbb(I=0)=\Pbb(A^c)=1/2\). Thus
\[
\mathcal L(I)=\frac12\delta_0+\frac12\mathcal N(0,1),
\]
which is not Gaussian. Random predictable integrands do not generally produce
Gaussian stochastic integrals.
\end{shortsolution}

\subsection{The accumulated gain process}

For a simple process, define
\[
M_t=\int_0^t\phi_s\dd W_s.
\]
On \((t_i,t_{i+1}]\),
\[
M_t=M_{t_i}+\xi_i(W_t-W_{t_i}).
\]
Thus the gain already accumulated is carried forward, while the current
position multiplies the fresh Brownian move. This interpretation survives the
\(L^2\) extension below.

% ============================================================
\section{The \texorpdfstring{\(\Ito\)}{Ito} Isometry and the
\texorpdfstring{\(L^2\)}{L2} Extension}
% ============================================================

\subsection{Orthogonality of fresh increments}

Let \(\Delta W_i=W_{t_{i+1}}-W_{t_i}\). Brownian increments on disjoint
intervals are independent, so
\[
\E[\Delta W_i\Delta W_j]=0\quad(i\neq j),
\qquad
\E[(\Delta W_i)^2]=t_{i+1}-t_i.
\]
For random predictable coefficients, the operational identity is
\[
\E[\xi_i\Delta W_i\mid\cF_{t_i}]
=\xi_i\E[\Delta W_i\mid\cF_{t_i}]=0.
\]
The later Brownian increment is fresh after conditioning on everything known
before it. This conditional centering is what removes every off-diagonal term
in the second-moment calculation.

\begin{theorem}[\(\Ito\) isometry for simple predictable processes]
If \(\phi\) is simple, predictable, and square-integrable, then
\[
\boxed{
\E\left[\left(\int_0^T\phi_t\dd W_t\right)^2\right]
=\E\int_0^T\phi_t^2\dd t.}
\]
Equivalently,
\[
\norm{\int_0^T\phi_t\dd W_t}_{L^2(\Omega)}
=\norm{\phi}_{L^2(\Omega\times[0,T])}.
\]
\end{theorem}

\begin{proof}
Write \(I=\sum_i\xi_i\Delta W_i\). Expanding the square gives
\[
\E[I^2]
=\sum_i\E[\xi_i^2(\Delta W_i)^2]
+2\sum_{i<j}\E[\xi_i\Delta W_i\,\xi_j\Delta W_j].
\]
If \(i<j\), then \(\xi_i\Delta W_i\xi_j\) is
\(\cF_{t_j}\)-measurable. Hence
\begin{align*}
\E[\xi_i\Delta W_i\,\xi_j\Delta W_j]
&=\E\!\left[
\xi_i\Delta W_i\xi_j
\E[\Delta W_j\mid\cF_{t_j}]
\right]\\
&=0.
\end{align*}
On the diagonal, \(\Delta W_i\) is independent of \(\cF_{t_i}\), so
\[
\E[\xi_i^2(\Delta W_i)^2]
=\E[\xi_i^2](t_{i+1}-t_i).
\]
Therefore
\[
\E[I^2]
=\sum_i\E[\xi_i^2](t_{i+1}-t_i)
=\E\int_0^T\phi_t^2\dd t.
\]
\end{proof}

\begin{keyidea}{Geometric interpretation}
The stochastic integral is an isometry: it transforms an exposure profile in
\(L^2(\Omega\times[0,T])\) into a terminal random gain in \(L^2(\Omega)\)
without changing its squared norm. Integrand energy becomes gain variance.
\end{keyidea}

\subsection{Mean, covariance, and risk bounds}

The same conditioning argument gives
\[
\E\left[\int_0^T\phi_t\dd W_t\right]=0.
\]
Applying the isometry to \(\phi+\psi\) and \(\phi-\psi\), or expanding
directly, gives the polarization identity.

\begin{proposition}[Covariance identity]
For square-integrable predictable \(\phi\) and \(\psi\),
\[
\boxed{
\E\left[
\left(\int_0^T\phi_t\dd W_t\right)
\left(\int_0^T\psi_t\dd W_t\right)
\right]
=\E\int_0^T\phi_t\psi_t\dd t.}
\]
\end{proposition}

Define the pathwise exposure and its expected energy by
\[
\mathcal E_T(\phi)=\int_0^T\phi_t^2\dd t,
\qquad
\E[\mathcal E_T(\phi)]=\E\int_0^T\phi_t^2\dd t.
\]
The isometry implies:
\begin{itemize}
  \item multiplying the position by \(c\) multiplies gain variance by \(c^2\);
  \item holding a fixed position twice as long doubles its variance contribution;
  \item moving exposure across time may preserve terminal variance while changing
        the whole gain path.
\end{itemize}
If \(|\phi_t|\leq L\), then
\[
\Var\left(\int_0^T\phi_t\dd W_t\right)
=\E\int_0^T\phi_t^2\dd t\leq L^2T.
\]

\begin{studentexercise}{same terminal risk, different timing}
On \([0,1]\), compare
\[
\phi_t=1,
\qquad
\psi_t=\sqrt2\,\1_{(0,1/2]}(t).
\]
Compute their energies and stochastic integrals. Compare their terminal laws
and their gain processes before time 1.
\end{studentexercise}

\begin{shortsolution}
Both strategies have unit energy:
\[
\int_0^1\phi_t^2\dd t=1,
\qquad
\int_0^1\psi_t^2\dd t=2\cdot\frac12=1.
\]
Moreover,
\[
\int_0^1\phi_t\dd W_t=W_1,
\qquad
\int_0^1\psi_t\dd W_t=\sqrt2\,W_{1/2}.
\]
Both terminal gains are \(\mathcal N(0,1)\). Their paths are not identical:
the second gain stops changing after time \(1/2\), while the first remains
exposed until time 1.
\end{shortsolution}

\subsection{Completion in \texorpdfstring{\(L^2\)}{L2}}

Let \(\phi\) be predictable with
\[
\E\int_0^T\phi_t^2\dd t<\infty.
\]
Choose simple predictable processes \(\phi^n\) such that
\[
\E\int_0^T|\phi_t^n-\phi_t|^2\dd t\longrightarrow0.
\]
The isometry gives
\begin{align*}
\E\left[
\left|
\int_0^T\phi_t^n\dd W_t-
\int_0^T\phi_t^m\dd W_t
\right|^2
\right]
&=\E\int_0^T|\phi_t^n-\phi_t^m|^2\dd t.
\end{align*}
The right-hand side tends to zero, so the sequence of integrals is Cauchy in
the complete space \(L^2(\Omega)\).

\begin{definition}[General \(\Ito\) integral]
For a square-integrable predictable process \(\phi\), define
\[
\boxed{
\int_0^T\phi_t\dd W_t
:=L^2\text{-}\lim_{n\to\infty}
\int_0^T\phi_t^n\dd W_t.}
\]
\end{definition}

The definition does not depend on the approximating sequence. Indeed, if
\(\psi^n\to\phi\) as well, the isometry yields
\[
\E\left[
\left|
\int_0^T\phi_t^n\dd W_t-
\int_0^T\psi_t^n\dd W_t
\right|^2
\right]
=\E\int_0^T|\phi_t^n-\psi_t^n|^2\dd t\longrightarrow0.
\]
Linearity, zero mean, the isometry, and the covariance identity all pass to
the limit.

\begin{keyidea}{The construction in one sentence}
Approximate the predictable exposure, integrate the simple approximations, and
take the \(L^2\) limit. The isometry simultaneously proves convergence,
uniqueness, and stability.
\end{keyidea}

% ============================================================
\section{The Integral as a Martingale Gain Process}
% ============================================================

For a predictable process \(\phi\) satisfying
\[
\E\int_0^T\phi_s^2\dd s<\infty,
\]
define the accumulated gain
\[
M_t=\int_0^t\phi_s\dd W_s,
\qquad 0\leq t\leq T.
\]

\subsection{The integrability audit}

Before calling \(M\) a martingale, three points must be checked:
\begin{enumerate}
  \item \(M_t\) is \(\cF_t\)-measurable;
  \item \(M_t\in L^1\) for every \(t\);
  \item \(\E[M_t\mid\cF_s]=M_s\) for \(s\leq t\).
\end{enumerate}
Adaptedness follows first for simple integrands and then by approximation. The
second point is exactly where the expected finite-energy condition matters.
The isometry and Cauchy--Schwarz give
\begin{align*}
\E|M_t|
&\leq \bigl(\E[M_t^2]\bigr)^{1/2}\\
&=\left(\E\int_0^t\phi_s^2\dd s\right)^{1/2}<\infty.
\end{align*}
Thus \(M_t\in L^2\subset L^1\). In particular, the conditional expectation
\(\E[M_t\mid\cF_s]\) is well defined.

\begin{warningbox}{True martingale versus local martingale}
There are two distinct assumptions:
\[
\boxed{\E\int_0^T\phi_s^2\dd s<\infty}
\qquad\text{and}\qquad
\boxed{\int_0^T\phi_s^2\dd s<\infty\quad\text{a.s.}}
\]
The first gives the square-integrable, hence true, martingale used throughout
this lecture. If only the second holds locally on every finite horizon, the
stochastic integral is in general only a continuous local martingale. A local
martingale need not be integrable at every time and need not preserve its
expectation. One must never silently replace the expected \(L^2\) condition by
pathwise finiteness.
\end{warningbox}

\subsection{Martingale property and proof}

\begin{theorem}[Square-integrable martingale property]
If \(\phi\) is predictable and
\(\E\int_0^T\phi_s^2\dd s<\infty\), then
\[
M_t=\int_0^t\phi_s\dd W_s
\]
has a continuous adapted version, \(M_t\in L^2\), and
\[
\boxed{\E[M_t\mid\cF_s]=M_s,
\qquad 0\leq s\leq t\leq T.}
\]
Equivalently,
\[
\E\left[\int_s^t\phi_u\dd W_u\middle|\cF_s\right]=0.
\]
\end{theorem}

\begin{proof}
First suppose that \(\phi\) is simple and use a grid containing \(s\) and
\(t\). Splitting past and future gains gives
\[
M_t=M_s+\sum_{t_i\geq s}\xi_i\Delta W_i.
\]
For each future term,
\[
\E[\xi_i\Delta W_i\mid\cF_{t_i}]
=\xi_i\E[\Delta W_i\mid\cF_{t_i}]=0.
\]
Successive applications of the tower property yield
\[
\E[M_t-M_s\mid\cF_s]=0.
\]

For a general integrand, choose simple \(\phi^n\to\phi\) in
\(L^2(\Omega\times[0,T])\), and let
\(M_t^n=\int_0^t\phi_u^n\dd W_u\). The isometry gives
\(M_t^n\to M_t\) in \(L^2\), hence in \(L^1\). Conditional expectation is an
\(L^1\)-contraction:
\[
\E\left|
\E[M_t^n-M_t\mid\cF_s]
\right|
\leq\E|M_t^n-M_t|\longrightarrow0.
\]
Therefore one may pass to the limit in
\(\E[M_t^n\mid\cF_s]=M_s^n\), proving the claim. The continuous version follows
from the standard construction of the Brownian stochastic integral.
\end{proof}

The proof contains the reusable mechanism of stochastic integration:
\[
\text{predictable coefficient}\times
\text{fresh centered increment}
\quad\Longrightarrow\quad
\text{zero conditional mean}.
\]

\subsection{Conditional variance and state-dependent risk}

The conditional version of the isometry is
\[
\boxed{
\E\left[
\left(\int_s^t\phi_u\dd W_u\right)^2
\middle|\cF_s
\right]
=\E\left[
\int_s^t\phi_u^2\dd u
\middle|\cF_s
\right].}
\]
Because the future integral has conditional mean zero, this identity is also
its conditional variance. The mean forecast is unchanged, while the amount of
risk still ahead may depend strongly on the current state.

\begin{studentexercise}{a state-dependent position}
On \([0,2]\), define
\[
\phi_t=\1_{(0,1]}(t)
+\bigl(1+\1_{\{W_1>0\}}\bigr)\1_{(1,2]}(t),
\qquad
M_t=\int_0^t\phi_s\dd W_s.
\]
Compute \(M_1\), \(\E[M_2\mid\cF_1]\), and
\(\Var(M_2-M_1\mid\cF_1)\).
\end{studentexercise}

\begin{shortsolution}
Up to time 1, the position is one, so \(M_1=W_1\). Since \(M\) is a
martingale,
\[
\E[M_2\mid\cF_1]=M_1=W_1.
\]
Over \((1,2]\), the position is already known at time 1. Therefore
\begin{align*}
\Var(M_2-M_1\mid\cF_1)
&=\bigl(1+\1_{\{W_1>0\}}\bigr)^2\\
&=\begin{cases}
1,&W_1\leq0,\\
4,&W_1>0.
\end{cases}
\end{align*}
Adaptation changes conditional variance without creating conditional drift.
\end{shortsolution}

\subsection{Quadratic variation as a random clock}

\begin{proposition}[Quadratic variation of a stochastic integral]
For
\[
M_t=\int_0^t\phi_s\dd W_s,
\]
one has
\[
\boxed{[M]_t=\int_0^t\phi_s^2\dd s.}
\]
\end{proposition}

For a simple strategy, the formula is already visible from
\[
\sum_i(\Delta M_i)^2
=\sum_i\xi_i^2(\Delta W_i)^2
\approx\sum_i\xi_i^2\Delta t_i.
\]
The process \(M\) therefore runs on the exposure clock generated by
\(\phi^2\). Calendar time advances uniformly, while the stochastic gain
accumulates variation faster when the absolute position is larger.

\subsection{Stopping an exposure}

Let \(\tau\) be a stopping time. The predictable strategy
\(\phi_s=\1_{\{s\leq\tau\}}\) holds one unit until \(\tau\) and zero afterward.
For every finite \(t\),
\[
\boxed{
\int_0^t\1_{\{s\leq\tau\}}\dd W_s=W_{t\wedge\tau}.}
\]
Crucially, its integrability is automatic on a bounded horizon:
\[
\E\int_0^t\1_{\{s\leq\tau\}}^2\dd s
=\E[t\wedge\tau]\leq t<\infty.
\]
The martingale property and isometry may therefore be applied legitimately:
\[
\boxed{
\E[W_{t\wedge\tau}]=0,
\qquad
\E[W_{t\wedge\tau}^2]=\E[t\wedge\tau].}
\]

\begin{boardexample}{stopping at an upper barrier}
Let
\[
\tau_a=\inf\{t\geq0:W_t=a\},
\qquad a>0.
\]
The stopped exposure is
\[
M_t=\int_0^t\1_{\{s\leq\tau_a\}}\dd W_s
=W_{t\wedge\tau_a}.
\]
Without using the density of \(\tau_a\),
\[
\E[M_t]=0,
\qquad
\E[M_t^2]=\E[t\wedge\tau_a].
\]
The first statement comes from the true-martingale property; the second comes
from the isometry.
\end{boardexample}

\begin{keyidea}{What predictability and integrability buy us}
For every predictable \(\phi\) satisfying
\(\E\int_0^T\phi_t^2\dd t<\infty\), the gain process is a
square-integrable martingale and
\[
\E\left[\int_0^T\phi_t\dd W_t\right]=0,
\qquad
\E\left[\left(\int_0^T\phi_t\dd W_t\right)^2\right]
=\E\int_0^T\phi_t^2\dd t.
\]
The distribution may nevertheless be Gaussian, a mixture, stopped, skewed, or
state-dependent. Zero conditional drift does not mean zero risk.
\end{keyidea}

% ============================================================
\section{Deterministic Integrands and Gaussian Applications}
% ============================================================

\subsection{Gaussianity and the \texorpdfstring{\(L^2\)}{L2} inner product}

When the integrand is deterministic, the stochastic integral has additional
structure.

\begin{theorem}[Distribution for a deterministic integrand]
If \(f\in L^2(0,T)\) is deterministic, then
\[
I_f=\int_0^Tf(t)\dd W_t
\sim\mathcal N\left(0,\int_0^Tf(t)^2\dd t\right).
\]
\end{theorem}

\begin{proof}
For a deterministic step function, the integral is a linear combination of
independent Gaussian increments and is therefore Gaussian. Choose deterministic
step functions \(f_n\to f\) in \(L^2(0,T)\). The isometry gives
\[
\int_0^Tf_n(t)\dd W_t\longrightarrow
\int_0^Tf(t)\dd W_t
\quad\text{in }L^2(\Omega),
\]
and the variances converge to \(\int_0^Tf(t)^2\dd t\). The characteristic
functions consequently converge to that of the stated centered Gaussian law.
\end{proof}

The same reasoning for linear combinations gives joint Gaussianity.

\begin{proposition}[Covariance and independence]
For deterministic \(f,g\in L^2(0,T)\),
\[
\Cov\left(\int_0^Tf\dd W,\int_0^Tg\dd W\right)
=\int_0^Tf(t)g(t)\dd t.
\]
The pair is jointly Gaussian. Therefore
\[
\int_0^Tf(t)g(t)\dd t=0
\quad\Longrightarrow\quad
\int_0^Tf\dd W\ \text{and}\ \int_0^Tg\dd W
\text{ are independent}.
\]
\end{proposition}

The map
\[
f\longmapsto\int_0^Tf\dd W
\]
turns the deterministic \(L^2(0,T)\) inner product into covariance. Ordinary
orthogonality of exposure profiles becomes independence of Gaussian risks.

\begin{boardexample}{a linearly increasing exposure}
Take \(f(t)=t\). Then
\[
I_T=\int_0^Tt\dd W_t
\]
is centered Gaussian with
\[
\Var(I_T)=\int_0^Tt^2\dd t=\frac{T^3}{3}.
\]
Thus
\[
\boxed{I_T\sim\mathcal N\left(0,\frac{T^3}{3}\right).}
\]
The standard deviation scales like \(T^{3/2}\), which is consistent with a
position of order \(T\) multiplied by Brownian noise of order \(T^{1/2}\).
\end{boardexample}

\subsection{Conditioning as orthogonal projection}

Let
\[
I_T=\int_0^Tt\dd W_t.
\]
Since
\[
W_T=\int_0^T1\dd W_t,
\]
the covariance identity gives
\[
\Cov(I_T,W_T)=\int_0^Tt\dd t=\frac{T^2}{2}.
\]
For a centered jointly Gaussian pair \((X,Y)\),
\[
\E[X\mid Y]=\frac{\Cov(X,Y)}{\Var(Y)}Y,
\qquad
\Var(X\mid Y)=\Var(X)-\frac{\Cov(X,Y)^2}{\Var(Y)}.
\]
Using \(\Var(W_T)=T\), we obtain
\[
\boxed{\E[I_T\mid W_T]=\frac{T}{2}W_T}
\]
and
\[
\boxed{
\Var(I_T\mid W_T)
=\frac{T^3}{3}-\frac{(T^2/2)^2}{T}
=\frac{T^3}{12}.}
\]

Geometrically, conditioning on \(W_T\) reveals the component of
\(t\mapsto t\) parallel to the constant function \(1\). Indeed,
\[
t=\frac{T}{2}+\left(t-\frac{T}{2}\right),
\qquad
\int_0^T\left(t-\frac{T}{2}\right)\dd t=0.
\]

\begin{studentexercise}{construct a Gaussian risk independent of \(W_T\)}
Find a nonzero affine function \(g(t)=t-c\) such that
\[
Y=\int_0^Tg(t)\dd W_t
\]
is independent of \(W_T\). Compute \(\Var(Y)\).
\end{studentexercise}

\begin{shortsolution}
Because \(W_T=\int_0^T1\dd W_t\), independence is equivalent to
\[
0=\Cov(Y,W_T)=\int_0^Tg(t)\dd t.
\]
For \(g(t)=t-c\), this gives \(T^2/2-cT=0\), hence \(c=T/2\). Therefore
\[
g(t)=t-\frac{T}{2},
\qquad
\Var(Y)=\int_0^T\left(t-\frac{T}{2}\right)^2\dd t
=\frac{T^3}{12}.
\]
The pair is jointly Gaussian and uncorrelated, so it is independent.
\end{shortsolution}

\subsection{Deterministic exposure as a time change}

For deterministic \(f\), define
\[
M_t=\int_0^tf(s)\dd W_s,
\qquad
q(t)=\int_0^tf(s)^2\dd s.
\]
Then \(M\) is centered Gaussian. For \(s<t\),
\[
M_t-M_s=\int_s^tf(u)\dd W_u
\]
depends only on Brownian increments over \((s,t]\), so disjoint increments of
\(M\) are independent. Moreover,
\[
\Var(M_t-M_s)=q(t)-q(s).
\]
Consequently,
\[
\boxed{(M_t)_{t\geq0}\overset{\mathrm{law}}=
(B_{q(t)})_{t\geq0}}
\]
for a Brownian motion \(B\). For \(f(t)=t\), the stochastic clock is
\(q(t)=t^3/3\).

This time-change statement is special to deterministic integrands. For a
random integrand, \(\int_0^t\phi_s^2\dd s\) is a random clock, but independent
increments and Gaussianity generally disappear.

\subsection{Integration by parts with a deterministic integrand}

The useful integration-by-parts formula in this lecture does not require
\(\Ito\)'s formula.

\begin{proposition}[Deterministic integration by parts]
Let \(f:[0,T]\to\R\) be deterministic and continuously differentiable. More
generally, assume that \(f\) is absolutely continuous with
\(f'\in L^2(0,T)\). Then
\[
f(T)W_T-f(0)W_0
=\int_0^Tf(t)\dd W_t+\int_0^TW_tf'(t)\dd t.
\]
Since \(W_0=0\),
\[
\boxed{
\int_0^Tf(t)\dd W_t
=f(T)W_T-\int_0^TW_tf'(t)\dd t.}
\]
\end{proposition}

\begin{proof}
The assumption \(f'\in L^2(0,T)\) implies \(f'\in L^1(0,T)\) on the finite
interval. Thus \(f\) has finite variation and belongs to \(L^2(0,T)\). On a
partition, discrete summation by parts gives
\[
f(T)W_T-f(0)W_0
=\sum_i f(t_i)\Delta W_i
+\sum_iW_{t_{i+1}}\bigl(f(t_{i+1})-f(t_i)\bigr).
\]
As the mesh tends to zero, the first sum converges in \(L^2\) to
\(\int_0^Tf(t)\dd W_t\). Since \(W\) is continuous and \(f\) has finite
variation, the second converges pathwise to
\[
\int_0^TW_t\dd f(t)=\int_0^TW_tf'(t)\dd t.
\]
This proves the formula directly from summation by parts.
\end{proof}

\begin{remark}[Why there is no correction term]
The deterministic path \(f\) has finite variation, so its quadratic
covariation with \(W\) is zero. No \(\Ito\) formula is needed for this result.
\end{remark}

\begin{boardexample}{the Brownian area}
Set
\[
A_T=\int_0^TW_t\dd t.
\]
Apply deterministic integration by parts with \(f(t)=t\):
\[
A_T=TW_T-\int_0^Tt\dd W_t.
\]
The right-hand side is a linear combination of jointly Gaussian stochastic
integrals, so \(A_T\) is centered Gaussian. Using
\[
\Var(W_T)=T,
\quad
\Var\left(\int_0^Tt\dd W_t\right)=\frac{T^3}{3},
\quad
\Cov\left(W_T,\int_0^Tt\dd W_t\right)=\frac{T^2}{2},
\]
we find
\begin{align*}
\Var(A_T)
&=T^2\Var(W_T)
+\Var\left(\int_0^Tt\dd W_t\right)
-2T\Cov\left(W_T,\int_0^Tt\dd W_t\right)\\
&=T^3+\frac{T^3}{3}-T^3
=\frac{T^3}{3}.
\end{align*}
Therefore
\[
\boxed{A_T\sim\mathcal N\left(0,\frac{T^3}{3}\right).}
\]
\end{boardexample}

% ============================================================
\section{Left, Right, and Midpoint Conventions}
% ============================================================

\subsection{One integrand, three sums}

For a partition \(\pi\) and the integrand \(\phi_t=W_t\), define
\begin{align*}
L^\pi&=\sum_iW_{t_i}\Delta W_i,\\
M^\pi&=\sum_i\frac{W_{t_i}+W_{t_{i+1}}}{2}\Delta W_i,\\
R^\pi&=\sum_iW_{t_{i+1}}\Delta W_i.
\end{align*}
Their information content is different:
\begin{center}
\begin{tabular}{lll}
\toprule
Sum & Evaluation rule & Interpretation \\
\midrule
left & value known before the move & \(\Ito\) \\
midpoint & uses both endpoints & Stratonovich \\
right & value known after the move & anticipative here \\
\bottomrule
\end{tabular}
\end{center}

The exact identity
\[
W_{t_{i+1}}^2-W_{t_i}^2
=2W_{t_i}\Delta W_i+(\Delta W_i)^2
\]
telescopes to
\[
W_T^2=2L^\pi+\sum_i(\Delta W_i)^2.
\]
Since Brownian quadratic variation converges to \(T\),
\[
\boxed{
\int_0^TW_t\dd W_t
=\lim_{|\pi|\to0}L^\pi
=\frac12(W_T^2-T).}
\]
This is the first nontrivial stochastic integral. The term \(-T/2\) is not
obtained from the general \(\Ito\) formula here; it follows directly from the
telescoping identity and quadratic variation.

Furthermore,
\[
R^\pi=L^\pi+\sum_i(\Delta W_i)^2,
\qquad
M^\pi=\frac12(L^\pi+R^\pi).
\]
Therefore the three limits are
\begin{center}
\begin{tabular}{lc}
\toprule
Convention & Limit \\
\midrule
left / \(\Ito\) & \(\frac12(W_T^2-T)\) \\
midpoint / Stratonovich & \(\frac12W_T^2\) \\
right & \(\frac12(W_T^2+T)\) \\
\bottomrule
\end{tabular}
\end{center}

\begin{studentexercise}{verify the isometry for \(\int W\dd W\)}
Let
\[
I_T=\int_0^TW_t\dd W_t=\frac12(W_T^2-T).
\]
Compute \(\E[I_T]\) and \(\E[I_T^2]\), using
\(\E[W_T^4]=3T^2\). Check the answer with the \(\Ito\) isometry.
\end{studentexercise}

\begin{shortsolution}
Since \(\E[W_T^2]=T\),
\[
\E[I_T]=\frac12(\E[W_T^2]-T)=0.
\]
Moreover,
\[
\E[I_T^2]
=\frac14\E[(W_T^2-T)^2]
=\frac14(3T^2-2T^2+T^2)
=\frac{T^2}{2}.
\]
The isometry gives exactly the same result:
\[
\E[I_T^2]
=\E\int_0^TW_t^2\dd t
=\int_0^T\E[W_t^2]\dd t
=\int_0^Tt\dd t
=\frac{T^2}{2}.
\]
\end{shortsolution}

\subsection{Why anticipation creates a fake arbitrage}

Suppose illegally that one could choose
\[
\phi_{t_i}=\operatorname{sign}(\Delta W_i)
\]
before reporting the gain. Then
\[
\sum_i\phi_{t_i}\Delta W_i=\sum_i|\Delta W_i|\geq0,
\]
and the gain is strictly positive almost surely on every nontrivial grid.
Under grid refinement, this sum even reflects the infinite total variation of
Brownian paths.

The apparent opportunity comes entirely from hindsight:
\(\phi_{t_i}\notin\cF_{t_i}\). Predictability is therefore both the
mathematical boundary of the \(\Ito\) integral and the economic boundary of a
legal self-financing strategy.

% ============================================================
\clearpage
\section{Simulation of Stochastic Integrals}
% ============================================================

\subsection{The left-point algorithm}

Let \(t_i=i\Delta t\), where \(\Delta t=T/n\). Generate independent standard
normal variables \(Z_0,\ldots,Z_{n-1}\) and set
\[
\Delta W_i=\sqrt{\Delta t}\,Z_i.
\]
Starting from \(W_0=0\) and \(I_0=0\), repeat for \(i=0,\ldots,n-1\):
\begin{enumerate}
  \item compute \(\phi_i\) from information available at \(t_i\);
  \item set \(I_{i+1}=I_i+\phi_i\Delta W_i\);
  \item set \(W_{t_{i+1}}=W_{t_i}+\Delta W_i\).
\end{enumerate}
Then
\[
I_n=\sum_{i=0}^{n-1}\phi_i\Delta W_i
\]
is the left-point approximation of \(\int_0^T\phi_t\dd W_t\).

\begin{warningbox}{The update order encodes the filtration}
The value \(\phi_i\) must be fully determined before \(\Delta W_i\) is used.
If it is computed from \(W_{t_{i+1}}\), or from the newly generated increment,
the scheme is a right-point or otherwise anticipative sum rather than an
\(\Ito\) approximation.
\end{warningbox}

For a target integrand \(\phi\), define the left-step approximation
\[
\phi_t^{(n)}=\sum_{i=0}^{n-1}\phi_{t_i}\1_{(t_i,t_{i+1}]}(t).
\]
Whenever \(\phi^{(n)}\to\phi\) in
\(L^2(\Omega\times[0,T])\), the isometry gives the exact mean-square error
identity
\[
\boxed{
\E\left[
\left|
\int_0^T\phi_t\dd W_t-
\sum_{i=0}^{n-1}\phi_{t_i}\Delta W_i
\right|^2
\right]
=\E\int_0^T|\phi_t-\phi_t^{(n)}|^2\dd t.}
\]
Thus the numerical approximation is controlled by approximation of the
integrand in the same \(L^2\) norm used in the theoretical construction.

\subsection{An exact benchmark}

For
\[
I=\int_0^TW_t\dd W_t,
\qquad
I^{(n)}=\sum_{i=0}^{n-1}W_{t_i}\Delta W_i,
\]
the identity \(I=\frac12(W_T^2-T)\) gives
\[
I-I^{(n)}
=\frac12\sum_{i=0}^{n-1}\bigl((\Delta W_i)^2-\Delta t\bigr).
\]
For a uniform grid,
\[
\Var((\Delta W_i)^2)=2(\Delta t)^2.
\]
The centered summands are independent, so
\begin{align*}
\E[(I-I^{(n)})^2]
&=\frac14\sum_{i=0}^{n-1}2(\Delta t)^2\\
&=\frac12n\left(\frac{T}{n}\right)^2
=\boxed{\frac{T^2}{2n}}.
\end{align*}
The root-mean-square error is therefore
\[
\boxed{\operatorname{RMSE}(I^{(n)})=\frac{T}{\sqrt{2n}}.}
\]
This supplies a clean convergence-rate test for an implementation.

\subsection{Diagnostics and classic traps}

\begin{center}
\small
\begin{tabular}{p{0.25\linewidth}p{0.35\linewidth}p{0.30\linewidth}}
\toprule
Mistake & Consequence & Check or correction \\
\midrule
Use \(\Delta t\,Z_i\) for \(\Delta W_i\)
& Brownian variance is too small and has the wrong time scaling
& Verify \(\Delta W_i=\sqrt{\Delta t}\,Z_i\) and empirical variance \(\Delta t\) \\
\addlinespace
Compute \(\phi_i\) from \(W_{t_{i+1}}\)
& The scheme anticipates and can introduce an order-one bias
& Freeze \(\phi_i\) from the state at \(t_i\) \\
\addlinespace
Use fresh independent noise for a jointly driven integrand
& The covariance structure and joint law are changed
& Reuse the same Brownian increments whenever the model specifies the same driver \\
\addlinespace
Replace \(\Delta W_i\) by \(\Delta t\)
& A stochastic integral is replaced by an ordinary time integral
& Track \(\dd W\) and \(\dd t\) terms separately \\
\addlinespace
Check only the empirical mean
& Variance or dependence errors remain invisible
& Test the isometry, covariance, and several grid sizes \\
\bottomrule
\end{tabular}
\end{center}

A single plausible-looking path is not validation. Useful checks include:
\begin{itemize}
  \item sample mean close to zero when the hypotheses imply a true martingale;
  \item sample second moment close to
        \(\E\int_0^T\phi_t^2\dd t\);
  \item covariance checks when several integrals use the same Brownian driver;
  \item error decay under successive grid refinements;
  \item comparison with an exact identity when one is available.
\end{itemize}

\begin{keyidea}{Final checklist}
\begin{enumerate}
  \item The \(\Ito\) integral is the \(L^2\) limit of left-point,
        non-anticipative gains.
  \item Predictability controls information; expected finite energy controls
        integrability.
  \item Under \(\E\int_0^T\phi_t^2\dd t<\infty\), the gain process is a
        square-integrable martingale with quadratic variation
        \(\int_0^t\phi_s^2\dd s\).
  \item Deterministic integrands produce jointly Gaussian variables, and
        \(L^2\) inner products become covariances.
  \item Deterministic integration by parts follows from summation by parts and
        requires no \(\Ito\) formula.
  \item Brownian quadratic variation produces
        \(\int_0^TW_t\dd W_t=\frac12(W_T^2-T)\).
  \item A correct simulation respects the same filtration and \(L^2\)
        structure as the theoretical integral.
\end{enumerate}
\end{keyidea}

\end{document}
